\documentclass[11pt,a4paper]{article}
\usepackage[T1]{fontenc}
\usepackage[utf8]{inputenc}
\usepackage{amsmath,amssymb,amsthm}
\usepackage[margin=1in]{geometry}
\usepackage[colorlinks=true,linkcolor=blue,urlcolor=blue]{hyperref}
\theoremstyle{plain}
\newtheorem{theorem}{Theorem}
\newtheorem{lemma}[theorem]{Lemma}
\newtheorem{proposition}[theorem]{Proposition}
\newtheorem{corollary}[theorem]{Corollary}
\theoremstyle{definition}
\newtheorem{remark}[theorem]{Remark}

\title{The empirical recurrence for OEIS A253360, proved by transfer matrix}
\author{Adrian Perez Fontelles\\ \small Independent researcher}
\date{5 September 2026}
\begin{document}
\maketitle

\begin{abstract}
OEIS A253360 counts arrays of width $3$ over $\{0,1\}$ in which the largest entry on the main diagonal minus the smallest on the antidiagonal of every $2\times2$ and $3\times3$ subblock differs from the same number at the neighbouring subblock. The entry carries an
empirical recurrence of order $6$ contributed by R.~H.~Hardin. It is true, and it is
decidable rather than empirical: a block of $k$ rows and the block one row lower together
occupy $k+1$ consecutive rows, so every comparison the condition makes is settled inside
$4$ consecutive rows and the count is a walk count on $S=201$ states.
\end{abstract}

\noindent\small 2020 Mathematics Subject Classification. 05A15, 05B20, 15A18.\normalsize

\section{The conjecture}

OEIS A253360 is ``Number of (n+2) X (1+2) 0..1 arrays with every 2 X 2 and 3 X 3 subblock diagonal maximum minus antidiagonal minimum unequal to its neighbors horizontally and vertically.''. It has offset $1$ and begins
\[
152, 272, 560, 976, 1904, 3472, 6544, 12048,\ \dots
\]
The entry states, as a conjecture contributed by its author and never marked settled:
\begin{quote}
Empirical: a(n) = 4*a(n-2) - 3*a(n-4) + 4*a(n-6).
\end{quote}
As of the ``Last modified'' line on the live entry (Dec 11 2018 09:00:05, revision 8) this is still
recorded as empirical, and nothing on the entry records it as proved.

\section{What is being counted}

The arrays have $3$ columns and a growing number of rows, with entries in $\{0,1\}$.

For a $k\times k$ block $B$ of cells of the array, let
\[
f(B)=\max_{1\le t\le k}B_{t,t}\;-\;\min_{1\le t\le k}B_{t,k+1-t} ,
\]
which in words is the largest entry on the main diagonal minus the smallest on the antidiagonal. The condition is that $f$ takes different values on any two blocks that
are neighbours in one of the directions the entry names: writing $B(i,j)$ for the block whose
top-left cell is $(i,j)$,
\[
f\bigl(B(i,j)\bigr)\;\neq\;f\bigl(B(i+\delta_1,\,j+\delta_2)\bigr)
\qquad\text{for }(\delta_1,\delta_2)\in\Delta,
\]
whenever both blocks lie wholly inside the array, where
\[
\Delta=\{\,(0,1),\ (1,0)\,\}
\]
collects the steps for the directions the entry names, namely horizontal and vertical. The
requirement is imposed for each of the two sizes of block, $2\times2$ and $3\times3$.

Two remarks fix the reading. A neighbouring block is the block one step over, overlapping the
first, not the next disjoint block: the entry says ``its neighbors'', and the counts below
confirm it. And a comparison is imposed only when both blocks lie wholly inside the array, so a
block near an edge is compared in fewer directions and a block whose named neighbour falls off
the array is under no condition in that direction.

\section{A window of 4 rows}

A $k\times k$ block occupies $k$ consecutive rows, and each of its neighbours occupies either
the same $k$ rows or the $k$ rows one lower. So both blocks of any comparison lie inside $k+1$
consecutive rows, hence inside $4$ consecutive rows for every size the entry names. No
comparison reaches further than that, in either direction.

Take the array one row at a time and let the state be the last 3 rows written so far, with a
distinguished start standing for the array before its first row, so that a walk of $n$ steps is
an array of $n$ rows. Appending a row produces a window of $4$ rows, and the step is
allowed exactly when every comparison whose two blocks both lie inside that window holds.

Every comparison of the finished array is tested this way. Each one lies inside $4$
consecutive rows, and every $4$ consecutive rows of the array form the window of some
step. A comparison lying inside more than one window is tested more than once, which costs
nothing. Conversely no step imposes anything not asked by the definition, since a step tests
only comparisons whose two blocks lie inside the array.

Every state may end an array, so $\tau=\mathbf{1}$. A block sitting in the final rows whose
neighbour would lie below the last row is not compared in that direction at all --- the
neighbour is not part of the array --- so stopping never leaves a condition unchecked. For the
same reason the array of exactly $3$ rows is handled by the walks of length $3$ with no
further test: all of its comparisons already lie inside the windows of those steps.

Thus
\[
a(n)\;=\;\iota^{\!\top}M^{\,j}\tau ,\qquad S=201,
\]
for the transfer matrix $M$ on those $S$ states. The walk index $j$ and the entry's own $n$
differ by a constant, read off by matching the published terms rather than assumed. In
particular $a$ is $C$-finite.

\section{The proof}

\begin{lemma}\label{lem:crit}
Let $q(t)=t^{6}-\sum_i c_it^{6-i}$ be the characteristic polynomial of the
conjectured recurrence. The residual $a(n)-\sum_ic_ia(n-i)$ equals
$\iota^{\!\top}M^{\,j}q(M)\tau$ for the corresponding walk index $j$, so the recurrence
holds from some point on if and only if $\iota^{\!\top}M^{j}q(M)\tau=0$ for all large $j$,
and it is enough to examine $j<S$.
\end{lemma}

\begin{proof}
Factor $M^{j}$ out of $M^{j+6}-\sum_ic_iM^{j+6-i}$. For the bound, $M^{S}$
is an integer combination of $I,M,\dots,M^{S-1}$ by Cayley--Hamilton, so the numbers
$u_j=\iota^{\!\top}M^{j}q(M)\tau$ obey the monic recurrence given by the characteristic
polynomial of $M$; $S$ consecutive zeros force every later one, and the last nonzero $u_j$
pins the threshold exactly.
\end{proof}

\begin{theorem}
\[
a(n)\;=\;4\,a(n-2) - 3\,a(n-4) + 4\,a(n-6)
\]
for every $n>6$.
\end{theorem}

\begin{proof}
The vector $w=q(M)\tau$ was formed by $6$ matrix--vector products in exact integer
arithmetic and $\iota^{\!\top}M^{j}w$ evaluated for $j=0,1,\dots$, again exactly; those
integers vanish from the index corresponding to $n=6+1$ onwards and the run continued
until the working vector was identically zero, which settles every larger $j$ at once.
\end{proof}

\section{Verification}

Four checks, each independent of the algebra above.

First, the model reproduces all $28$ terms the entry publishes, exactly, in integer
arithmetic.

Second, the count was recomputed for the smallest arrays straight from the definition: every
array of width $3$ over $\{0,1\}$ with a few rows was written out cell by cell, every block of
each of the two sizes was located, its number computed, and the comparison made with each neighbouring block
that lies inside the array. No window, no state and no end vector are involved in that
computation, and the published terms came back.

Third, the conjectured recurrence was evaluated directly on the published terms, with no
matrices involved, and holds wherever the proved range and the data overlap.

Fourth, the annihilation test of Lemma~\ref{lem:crit} was carried out in exact integer
arithmetic on the whole vector rather than on a sampled prefix.

\begin{thebibliography}{9}
\bibitem{oeis} The OEIS Foundation, \emph{The On-Line Encyclopedia of Integer
Sequences}, \url{https://oeis.org}, 2026. Sequence A253360.
\bibitem{stanley} R.~P.~Stanley, \emph{Enumerative Combinatorics}, Volume 1, 2nd ed.,
Cambridge University Press, 2011. (Transfer-matrix method, Section 4.7.)
\bibitem{horn} R.~A.~Horn and C.~R.~Johnson, \emph{Matrix Analysis}, 2nd ed., Cambridge
University Press, 2013.
\end{thebibliography}

\end{document}
